% SpecLens · SANER 2028 投稿骨架（I5 起步，2026-11-15）
% 格式：正式投稿时换 IEEEtran 双栏（SANER 用 IEEE CNF 模板）；当前以 article 便于内容迭代。
% 数据来源：experiments/runs/litmus-*（deepseek/deepseek-flash, temperature 0, run_index 0）
%   与 litmus-minimax-*（第二模型）；所有表格数字须由 runner 产物重新生成后再入稿（勿手抄）。
\documentclass[11pt]{article}
\usepackage[margin=1in]{geometry}
\usepackage{booktabs}
\usepackage{graphicx}
\usepackage[hidelinks]{hyperref}

\title{SpecLens: Measuring and Improving LLM Code Review\\Against Unseen Specifications}
\author{Anonymous}
\date{}

\begin{document}
\maketitle

\begin{abstract}
% 论点主线（按 I2-I4 实证更新）：
% 1) 规范在场使幻觉误报清零（FP 0.539 -> 0.077, p=0.031）——「知识注入」有效但已知；
% 2) 全文 in-context 的有效性依赖「塞得下」：预算 236K->30K 检出 0.870->0.217（RQ4 失效模式）；
% 3) 「规则找代码」vs「代码找规则」：B2 检索 1.3% 成本 0.478 检出（覆盖由命中决定）；
% 4) 已知规则集下，编译静态检查器零成本追平全文 LLM 基线（B3 0.870/0.000）；
% 5) LLM 自审式规则验证会静默移除关键 MUST 规则（3/12），应作人工复核信号而非过滤器。
\end{abstract}

\section{Introduction}
% RQ1（测量）：LLM 对照未见规范的遵守率/误报率；RQ3（归因）：错误分类学；
% RQ4（成本/规模）：预算阶梯 + 多仓库；RQ2（方法）：验证管线增益——当前实证不支持
% 其作为主贡献（I2 判读 Go-RQ4），论文主线 = benchmark + measurement + 成本-效果。

\section{Related Work}
% 近邻（idea.md 定位修订 v2 表格）：LintCFG、RulePilot、GPTAid、AutoCommenter、
% Code Review Agent Benchmark、Grounded Agentic Review；商业：Copilot instructions/
% CodeRabbit/Macroscope。错位点：规则的可验证性与遵守的一致性。

\section{The SpecLens-Bench Benchmark}
% 4 规范快照（来源/许可/commit pin）；40 条 react19 标注（v0.2：TS-011 修订后 100% 人工确认）；
% 难度分级与规则类型学（T1-T5）；标注指南 v1.2；240 条扩样进行中（I5）。
% 新鲜度探测：deepseek-flash 对全部 4 规范 0.90-1.00（「未见」前提不成立的诚实报告）。

\section{Study Design}
% 方法臂：B0/B1/B2(RAG)/B3(static)/M-lite/M-verified(fail-open)；控制变量原则；
% 记账（usage.json）与统计（精确二项 McNemar + violation/compliant 分组）。

\section{Results}

\subsection{RQ1: Specification Adherence (Measurement)}
Without the specification (B0), the model detects 91.3\% of violations but flags
6 of 13 compliant files (FP 0.539); with the full specification in context (B1),
false positives drop to 0.077 (McNemar $p=0.031$) while detection is statistically
unchanged. The mechanism, confirmed by our error taxonomy, is that zero-shot review
substitutes plausible software-engineering standards (error handling, caching,
concurrency hygiene) for the actual specification---a failure class we call
\emph{hallucinated standard}, which disappears entirely once the specification is
present. On the second model (MiniMax-M3.1-Flash) the same class reproduces
(FP 0.615 $
ightarrow$ 0.231), but suppression is weaker, showing that the
adherence gap is heterogeneous across models.

\subsection{RQ2: Does Rule Verification Help?}
Our verification pipeline (grounding, executability self-review, conflict check)
does not improve detection over M-lite: 0.870 vs 0.913 on the single-model pilot.
Root-cause analysis over archived per-rule gate reports shows the executability
gate falsely failed a MUST-level rule whose quote is verbatim from the snapshot
(grounding 1.0), silently removing it and costing 13 detection points; the same
gate passed the identical rule on a re-run at temperature 0. We therefore adopt a
fail-open policy---LLM self-review marks rules for human review instead of
removing them---and conclude that LLM self-review should serve as a review signal,
not a filter.

\subsection{RQ3: Error Taxonomy}
We classify every false positive and false negative of five configurations into
six classes. Beyond hallucinated standards, the notable classes are
\emph{rule inversion} (a second model inverts a prohibition into a requirement),
\emph{rule over-generalization} (per-rule review lacking document context flags
mandated caching as a risk), and \emph{soft-wording misses} (SHOULD-level
descriptive rules are systematically under-detected). One benchmark artifact---a
compliant case containing an out-of-scope real violation pattern, flagged
consistently by all five configurations---was resolved by revising the case,
illustrating how consistent cross-configuration disagreement localizes annotation
granularity problems.

\subsection{RQ4: Cost and Scale}
Full-specification review is effective only while the specification fits in
context: shrinking the budget from 236K to 30K characters collapses detection
from 0.870 to 0.217 while false positives stay low. On real repository gradients
(1k/10k/50k-line generated repositories with an embedded violation), B1 consumes
4.21M tokens for the 18-case suite versus 68K for rule-routed review---a
62$	imes$ gap at identical effectiveness---because routing resolves which files
a rule applies to and benign files cost nothing. A lexical-retrieval baseline
reaches 0.478 detection at 1.3\% of B1's cost, and a compiled static checker
matches B1's detection and false-positive profile at zero marginal cost,
delineating when knowledge organization beats context stuffing.

\section{Threats to Validity}
% 单模型已由第二模型对照补齐（MiniMax-M2.1，待 I5 写入）；规范熟悉度（probe 0.90-1.00）；
% 40 条规模；synthetic->handcrafted 复核流程；单文件 case 的生态效度。

\section{Conclusion}

\bibliographystyle{IEEEtran}
\bibliography{refs}
\end{document}
